% Curs 10 - Pagina 7 (6-XII-1994).
% Fragment pentru % Curs 10 - Pagina 7 (6-XII-1994).
% Fragment pentru % Curs 10 - Pagina 7 (6-XII-1994).
% Fragment pentru % Curs 10 - Pagina 7 (6-XII-1994).
% Fragment pentru \input{courses/pages/curs10-pagina7-ro}.
% Fara preambul si fara intrerupere fortata de pagina.
\bigskip

\subsection*{Curs 10 - Pagina 7}

\noindent $aPa^{-1}=f_a(P)$, unde $f_a$ este automorfism.
\[
\Longrightarrow\quad aPa^{-1}\leq G
\quad\text{şi}\quad |aPa^{-1}|=|P|.
\]

\bigskip
\noindent\fbox{\textbf{T. lui Sylow 2}}. Aceleaşi ipoteze ca la T. 1,
şi $H$ este $p$-subgrup al lui $G$, iar $P$ este $p$-subgrup Sylow al lui $G$.

\medskip
\noindent Atunci $(\exists)\,a\in G$ a.î.
\[
H\leq aPa^{-1}\in\operatorname{Syl}_p(G).
\]

\noindent Dacă $P,Q\in\operatorname{Syl}_p(G)$, atunci
$(\exists)\,a\in G$ a.î.
\[
aQa^{-1}=P.
\]

\bigskip
\noindent\fbox{\textbf{T. lui Sylow 3}}. Aceleaşi ipoteze ca la T. 1.
\[
|\operatorname{Syl}_p(G)|\mid |G|,
\]
\noindent şi $|\operatorname{Syl}_p(G)|$ este de forma
$p^k+1$\footnote{În manuscris exponentul pare să fie $k$,
ceea ce ar da $p^k+1$. Forma uzuală a teoremei este $kp+1$.
Am păstrat aici lectura vizibilă; este necesară verificarea originalului.},
cu
\[
|\operatorname{Syl}_p(G)|\equiv 1\pmod p.
\]

\bigskip
\noindent\underline{Aplicaţie a teoremelor Sylow în clasificarea grupurilor finite}.

\medskip
\noindent Dat $n\in\mathbb{N}^{*}$, să găsim o listă de grupuri de ordin $n$,
$G_1,\ldots,G_q$, a.î.:
\begin{itemize}
\item $(\forall)\,G$ grup de ordin $n$, $(\exists)\,i$ a.î. $G\cong G_i$;
\item $(\forall)\,i\neq j$, $G_i\not\cong G_j$.
\end{itemize}

\noindent $n=6:\quad (\mathbb{Z}_6,+,\hat 0),\quad(S_3,\circ,e).$

\bigskip
\noindent\fbox{\textbf{T.}} Dacă $G$ este grup a.î.
$|G|=n=pq$, $p>q$, $p,q$ prime, şi $q\nmid(p-1)$,
atunci
\[
G\cong\mathbb{Z}_{pq}.
\]

\medskip
\noindent De pildă, $n=35=7\cdot5$; atunci $G\cong\mathbb{Z}_{35}$.
% Exemplul este foarte prescurtat in manuscris.

\noindent Să considerăm cazul în care $q\mid(p-1)$.
\[
q=2,\qquad p=11.
\]

\bigskip
\noindent\underline{Lema 2}. Fie $G$ grup finit şi $H$ un $p$-subgrup al lui $G$.
Atunci
\[
[N_G(H):H]\equiv[G:H]\pmod p.
\]

\medskip
\noindent\underline{Dem. lema 2}: Iau
\[
A=\{aH\mid a\in G\}.
\]
\noindent $H\times A\longrightarrow A$, acţiune prin translaţie:
\[
(h,aH)\longmapsto(ha)H.
\]
% Demonstratia continua pe pagina urmatoare a manuscrisului.
.
% Fara preambul si fara intrerupere fortata de pagina.
\bigskip

\subsection*{Curs 10 - Pagina 7}

\noindent $aPa^{-1}=f_a(P)$, unde $f_a$ este automorfism.
\[
\Longrightarrow\quad aPa^{-1}\leq G
\quad\text{şi}\quad |aPa^{-1}|=|P|.
\]

\bigskip
\noindent\fbox{\textbf{T. lui Sylow 2}}. Aceleaşi ipoteze ca la T. 1,
şi $H$ este $p$-subgrup al lui $G$, iar $P$ este $p$-subgrup Sylow al lui $G$.

\medskip
\noindent Atunci $(\exists)\,a\in G$ a.î.
\[
H\leq aPa^{-1}\in\operatorname{Syl}_p(G).
\]

\noindent Dacă $P,Q\in\operatorname{Syl}_p(G)$, atunci
$(\exists)\,a\in G$ a.î.
\[
aQa^{-1}=P.
\]

\bigskip
\noindent\fbox{\textbf{T. lui Sylow 3}}. Aceleaşi ipoteze ca la T. 1.
\[
|\operatorname{Syl}_p(G)|\mid |G|,
\]
\noindent şi $|\operatorname{Syl}_p(G)|$ este de forma
$p^k+1$\footnote{În manuscris exponentul pare să fie $k$,
ceea ce ar da $p^k+1$. Forma uzuală a teoremei este $kp+1$.
Am păstrat aici lectura vizibilă; este necesară verificarea originalului.},
cu
\[
|\operatorname{Syl}_p(G)|\equiv 1\pmod p.
\]

\bigskip
\noindent\underline{Aplicaţie a teoremelor Sylow în clasificarea grupurilor finite}.

\medskip
\noindent Dat $n\in\mathbb{N}^{*}$, să găsim o listă de grupuri de ordin $n$,
$G_1,\ldots,G_q$, a.î.:
\begin{itemize}
\item $(\forall)\,G$ grup de ordin $n$, $(\exists)\,i$ a.î. $G\cong G_i$;
\item $(\forall)\,i\neq j$, $G_i\not\cong G_j$.
\end{itemize}

\noindent $n=6:\quad (\mathbb{Z}_6,+,\hat 0),\quad(S_3,\circ,e).$

\bigskip
\noindent\fbox{\textbf{T.}} Dacă $G$ este grup a.î.
$|G|=n=pq$, $p>q$, $p,q$ prime, şi $q\nmid(p-1)$,
atunci
\[
G\cong\mathbb{Z}_{pq}.
\]

\medskip
\noindent De pildă, $n=35=7\cdot5$; atunci $G\cong\mathbb{Z}_{35}$.
% Exemplul este foarte prescurtat in manuscris.

\noindent Să considerăm cazul în care $q\mid(p-1)$.
\[
q=2,\qquad p=11.
\]

\bigskip
\noindent\underline{Lema 2}. Fie $G$ grup finit şi $H$ un $p$-subgrup al lui $G$.
Atunci
\[
[N_G(H):H]\equiv[G:H]\pmod p.
\]

\medskip
\noindent\underline{Dem. lema 2}: Iau
\[
A=\{aH\mid a\in G\}.
\]
\noindent $H\times A\longrightarrow A$, acţiune prin translaţie:
\[
(h,aH)\longmapsto(ha)H.
\]
% Demonstratia continua pe pagina urmatoare a manuscrisului.
.
% Fara preambul si fara intrerupere fortata de pagina.
\bigskip

\subsection*{Curs 10 - Pagina 7}

\noindent $aPa^{-1}=f_a(P)$, unde $f_a$ este automorfism.
\[
\Longrightarrow\quad aPa^{-1}\leq G
\quad\text{şi}\quad |aPa^{-1}|=|P|.
\]

\bigskip
\noindent\fbox{\textbf{T. lui Sylow 2}}. Aceleaşi ipoteze ca la T. 1,
şi $H$ este $p$-subgrup al lui $G$, iar $P$ este $p$-subgrup Sylow al lui $G$.

\medskip
\noindent Atunci $(\exists)\,a\in G$ a.î.
\[
H\leq aPa^{-1}\in\operatorname{Syl}_p(G).
\]

\noindent Dacă $P,Q\in\operatorname{Syl}_p(G)$, atunci
$(\exists)\,a\in G$ a.î.
\[
aQa^{-1}=P.
\]

\bigskip
\noindent\fbox{\textbf{T. lui Sylow 3}}. Aceleaşi ipoteze ca la T. 1.
\[
|\operatorname{Syl}_p(G)|\mid |G|,
\]
\noindent şi $|\operatorname{Syl}_p(G)|$ este de forma
$p^k+1$\footnote{În manuscris exponentul pare să fie $k$,
ceea ce ar da $p^k+1$. Forma uzuală a teoremei este $kp+1$.
Am păstrat aici lectura vizibilă; este necesară verificarea originalului.},
cu
\[
|\operatorname{Syl}_p(G)|\equiv 1\pmod p.
\]

\bigskip
\noindent\underline{Aplicaţie a teoremelor Sylow în clasificarea grupurilor finite}.

\medskip
\noindent Dat $n\in\mathbb{N}^{*}$, să găsim o listă de grupuri de ordin $n$,
$G_1,\ldots,G_q$, a.î.:
\begin{itemize}
\item $(\forall)\,G$ grup de ordin $n$, $(\exists)\,i$ a.î. $G\cong G_i$;
\item $(\forall)\,i\neq j$, $G_i\not\cong G_j$.
\end{itemize}

\noindent $n=6:\quad (\mathbb{Z}_6,+,\hat 0),\quad(S_3,\circ,e).$

\bigskip
\noindent\fbox{\textbf{T.}} Dacă $G$ este grup a.î.
$|G|=n=pq$, $p>q$, $p,q$ prime, şi $q\nmid(p-1)$,
atunci
\[
G\cong\mathbb{Z}_{pq}.
\]

\medskip
\noindent De pildă, $n=35=7\cdot5$; atunci $G\cong\mathbb{Z}_{35}$.
% Exemplul este foarte prescurtat in manuscris.

\noindent Să considerăm cazul în care $q\mid(p-1)$.
\[
q=2,\qquad p=11.
\]

\bigskip
\noindent\underline{Lema 2}. Fie $G$ grup finit şi $H$ un $p$-subgrup al lui $G$.
Atunci
\[
[N_G(H):H]\equiv[G:H]\pmod p.
\]

\medskip
\noindent\underline{Dem. lema 2}: Iau
\[
A=\{aH\mid a\in G\}.
\]
\noindent $H\times A\longrightarrow A$, acţiune prin translaţie:
\[
(h,aH)\longmapsto(ha)H.
\]
% Demonstratia continua pe pagina urmatoare a manuscrisului.
.
% Fara preambul si fara intrerupere fortata de pagina.
\bigskip

\subsection*{Curs 10 - Pagina 7}

\noindent $aPa^{-1}=f_a(P)$, unde $f_a$ este automorfism.
\[
\Longrightarrow\quad aPa^{-1}\leq G
\quad\text{şi}\quad |aPa^{-1}|=|P|.
\]

\bigskip
\noindent\fbox{\textbf{T. lui Sylow 2}}. Aceleaşi ipoteze ca la T. 1,
şi $H$ este $p$-subgrup al lui $G$, iar $P$ este $p$-subgrup Sylow al lui $G$.

\medskip
\noindent Atunci $(\exists)\,a\in G$ a.î.
\[
H\leq aPa^{-1}\in\operatorname{Syl}_p(G).
\]

\noindent Dacă $P,Q\in\operatorname{Syl}_p(G)$, atunci
$(\exists)\,a\in G$ a.î.
\[
aQa^{-1}=P.
\]

\bigskip
\noindent\fbox{\textbf{T. lui Sylow 3}}. Aceleaşi ipoteze ca la T. 1.
\[
|\operatorname{Syl}_p(G)|\mid |G|,
\]
\noindent şi $|\operatorname{Syl}_p(G)|$ este de forma
$p^k+1$\footnote{În manuscris exponentul pare să fie $k$,
ceea ce ar da $p^k+1$. Forma uzuală a teoremei este $kp+1$.
Am păstrat aici lectura vizibilă; este necesară verificarea originalului.},
cu
\[
|\operatorname{Syl}_p(G)|\equiv 1\pmod p.
\]

\bigskip
\noindent\underline{Aplicaţie a teoremelor Sylow în clasificarea grupurilor finite}.

\medskip
\noindent Dat $n\in\mathbb{N}^{*}$, să găsim o listă de grupuri de ordin $n$,
$G_1,\ldots,G_q$, a.î.:
\begin{itemize}
\item $(\forall)\,G$ grup de ordin $n$, $(\exists)\,i$ a.î. $G\cong G_i$;
\item $(\forall)\,i\neq j$, $G_i\not\cong G_j$.
\end{itemize}

\noindent $n=6:\quad (\mathbb{Z}_6,+,\hat 0),\quad(S_3,\circ,e).$

\bigskip
\noindent\fbox{\textbf{T.}} Dacă $G$ este grup a.î.
$|G|=n=pq$, $p>q$, $p,q$ prime, şi $q\nmid(p-1)$,
atunci
\[
G\cong\mathbb{Z}_{pq}.
\]

\medskip
\noindent De pildă, $n=35=7\cdot5$; atunci $G\cong\mathbb{Z}_{35}$.
% Exemplul este foarte prescurtat in manuscris.

\noindent Să considerăm cazul în care $q\mid(p-1)$.
\[
q=2,\qquad p=11.
\]

\bigskip
\noindent\underline{Lema 2}. Fie $G$ grup finit şi $H$ un $p$-subgrup al lui $G$.
Atunci
\[
[N_G(H):H]\equiv[G:H]\pmod p.
\]

\medskip
\noindent\underline{Dem. lema 2}: Iau
\[
A=\{aH\mid a\in G\}.
\]
\noindent $H\times A\longrightarrow A$, acţiune prin translaţie:
\[
(h,aH)\longmapsto(ha)H.
\]
% Demonstratia continua pe pagina urmatoare a manuscrisului.
